\chapter{Formal Models of Doxonomic Systems}
\label{ch:formal-models}

\epigraph{A theory has only the alternative of being right or wrong. A model has a third possibility: it may be right, but irrelevant.}{Manfred Eigen}

\noindent
This chapter develops the mathematical core of doxonomics: stock-and-flow models, epidemiological diffusion, Bayesian updating under unequal exposure, cascade models, and differential equation frameworks for belief dynamics.

\section{Stock-and-Flow Models}
\label{sec:ch11-stock-flow}

\begin{definition}[Belief Stock and Narrative Flow]
\label{def:stock-flow}
\index{stock-and-flow model}
The \emph{belief stock} $\belief{i}(t) \in [0,1]$ is the fraction of the population holding belief $i$ at time $t$.
The \emph{narrative flow} $\phi_i(t) \in \R_{\geq 0}$ is the rate at which narrative content supporting $i$ is produced and circulated.
The dynamics are:
\begin{equation}
\label{eq:stock-flow}
  \frac{d\belief{i}}{dt} = \alpha \cdot \phi_i(t) \cdot (1 - \belief{i}(t)) - \delta \cdot \belief{i}(t) + \beta \cdot \belief{i}(t)(1 - \belief{i}(t))
\end{equation}
where $\alpha$ is the persuasion rate, $\delta$ is the natural decay rate, and $\beta$ is the social reinforcement rate (peer effects).
\end{definition}

\begin{proposition}[Steady States]
\label{prop:steady-states}
Equation~\eqref{eq:stock-flow} has at most three steady states.
Setting $d\belief{i}/dt = 0$, the trivial steady state $\belief{i}^* = 0$ exists when $\phi_i = 0$.
For $\phi_i > 0$ and $\beta > 0$, there exist two non-trivial steady states:
\[
  \belief{i}^* = \frac{(\alpha\phi_i + \beta - \delta) \pm \sqrt{(\alpha\phi_i + \beta - \delta)^2 + 4\beta\alpha\phi_i}}{2\beta}
\]
The system exhibits bistability when $\beta$ is sufficiently large relative to $\delta$, corresponding to the lock-in phenomenon of Chapter~\ref{ch:common-sense-capture}.
\end{proposition}

\section{Epidemiological Diffusion Models}
\label{sec:ch11-sir}

Belief spread shares structural features with disease transmission.

\begin{model}[SIR Model for Belief]
\label{mod:sir-belief}
\index{SIR model}
Partition the population into three compartments:
\begin{itemize}[nosep]
  \item $S(t)$: susceptible (have not encountered the narrative),
  \item $I(t)$: infected (actively hold and transmit the belief),
  \item $R(t)$: recovered (have encountered and rejected the belief, or lost interest).
\end{itemize}
The dynamics are:
\begin{align}
  \frac{dS}{dt} &= -\beta S I - \alpha \phi(t) S \label{eq:sir-s}\\
  \frac{dI}{dt} &= \beta S I + \alpha \phi(t) S - \gamma I \label{eq:sir-i}\\
  \frac{dR}{dt} &= \gamma I \label{eq:sir-r}
\end{align}
where $\beta$ is the peer transmission rate, $\alpha\phi(t)$ is the institutional transmission rate (from narrative infrastructure), and $\gamma$ is the recovery rate (skepticism, counter-narrative exposure).
\end{model}

\begin{definition}[Basic Reproduction Number]
\label{def:R0}
\index{basic reproduction number}
The \emph{doxonomic reproduction number} is
\[
  R_0 = \frac{\beta + \alpha\phi}{\gamma}
\]
If $R_0 > 1$, the belief can spread epidemically.
If $R_0 < 1$, it dies out without sustained institutional support ($\phi > 0$).
\end{definition}

A crucial difference from epidemiology: institutions can maintain $\phi > 0$ indefinitely, sustaining beliefs that would otherwise die out.
This is the formal basis of the claim that some beliefs persist not because they are persuasive but because they are subsidized.

\section{Bayesian Updating Under Unequal Exposure}
\label{sec:ch11-bayesian}

\begin{model}[Asymmetric Information Environment]
\label{mod:asymmetric-bayes}
\index{Bayesian updating!asymmetric}
Agent $j$ observes a sequence of signals $\{s_t^{(j)}\}_{t=1}^{T}$ where each $s_t^{(j)}$ is drawn from
\[
  s_t^{(j)} \sim \begin{cases} \nu^+_j & \text{with probability } q_j \\ \nu^-_j & \text{with probability } 1 - q_j \end{cases}
\]
and $q_j$ depends on $j$'s media environment.
Agents in channel $k$ have $q_j = q_k$, determined by the channel's editorial stance.
After $T$ signals, the posterior is
\[
  \pi_j^{(T)}(p) = \frac{\pi_j^{(0)}(p) \cdot \prod_{t=1}^T L(s_t^{(j)} \mid p)}{\pi_j^{(0)}(p) \cdot \prod_{t=1}^T L(s_t^{(j)} \mid p) + (1-\pi_j^{(0)}(p)) \cdot \prod_{t=1}^T L(s_t^{(j)} \mid \neg p)}
\]
\end{model}

\begin{theorem}[Polarization Under Asymmetric Exposure]
\label{thm:polarization}
\index{polarization}
If two groups $A$ and $B$ have $q_A > 1/2$ and $q_B < 1/2$ (their channels systematically favor opposite conclusions), then as $T \to \infty$:
\[
  \pi_A^{(T)}(p) \to 1 \quad \text{and} \quad \pi_B^{(T)}(p) \to 0
\]
almost surely, even if both groups are perfectly rational Bayesian updaters.
\end{theorem}

This theorem is the formal foundation for the claim that doxonomic polarization does not require irrationality---only asymmetric evidence supply.

\section{Cascade Models}
\label{sec:ch11-cascades}

\begin{model}[Information Cascade]
\label{mod:cascade}
\index{information cascade}
Following \textcite{banerjee1992} and \textcite{bikhchandani1992}, agents act sequentially, observing predecessors' actions but not their private signals.
An \emph{information cascade} occurs when agent $j$'s optimal action is independent of $j$'s private signal---the public evidence overwhelms private information.
In doxonomic terms, a cascade occurs when the accumulated weight of institutional narrative is so strong that individual experience becomes irrelevant to expressed belief.
\end{model}

\section{Differential Equation Framework}
\label{sec:ch11-ode}

We can model competing beliefs as a dynamical system.
Let $\belief{1}(t)$ and $\belief{2}(t) = 1 - \belief{1}(t)$ be two competing beliefs.

\begin{model}[Lotka-Volterra Belief Competition]
\label{mod:lotka-volterra}
\index{Lotka-Volterra model}
\begin{align}
  \frac{d\belief{1}}{dt} &= \belief{1}(\alpha_1 \phi_1 - \gamma_1 - \beta_{12}\belief{2}) \\
  \frac{d\belief{2}}{dt} &= \belief{2}(\alpha_2 \phi_2 - \gamma_2 - \beta_{21}\belief{1})
\end{align}
where $\phi_k$ is the narrative flow for belief $k$, $\gamma_k$ is the intrinsic decay, and $\beta_{ij}$ is the competitive displacement rate.
The funding asymmetry $\phi_1 \gg \phi_2$ can ensure dominance of belief 1 regardless of intrinsic merit.
\end{model}

\section{Notes and References}
\label{sec:ch11-notes}

Stock-and-flow models adapt systems dynamics \autocite{jackson2008}.
The SIR framework for information is developed in \textcite{watts2002}.
Bayesian polarization under asymmetric information is treated in the economics of media bias \autocite{gentzkow2010}.
Cascade models originate with \textcite{banerjee1992} and \textcite{bikhchandani1992}.
Lotka-Volterra models for cultural competition appear in cultural evolution literature \autocite{nowak2006}.

\section{Exercises}

\begin{exercise}
For the stock-and-flow model~\eqref{eq:stock-flow} with $\alpha = 0.3$, $\delta = 0.1$, $\beta = 0.5$, find the minimum narrative flow $\phi^*$ needed for a stable high-adoption equilibrium to exist.
\end{exercise}

\begin{exercise}
Compute $R_0$ for a belief spread through social media ($\beta = 0.05$, $\gamma = 0.2$) with institutional support $\alpha\phi = 0.1$.
Can the belief sustain itself without institutional support?
\end{exercise}

\begin{exercise}
Simulate the Bayesian polarization theorem: generate $T = 100$ signals for two groups with $q_A = 0.7$ and $q_B = 0.3$, starting from the same prior $\pi^{(0)} = 0.5$.
Plot the posterior trajectories.
\end{exercise}

\begin{exercise}
In the Lotka-Volterra model, find the conditions under which belief 2 survives despite $\phi_1 > \phi_2$.
Interpret economically: what makes a counter-narrative resilient?
\end{exercise}

\begin{exercise}
Extend the SIR model to include a ``vaccinated'' compartment $V(t)$ representing people who have received counter-narrative inoculation.
Write the modified equations and find the new $R_0$.
\end{exercise}
